\section{Connection to online learning theory}
\label{sec:ol}

\subsection{CFR as FTRL}

CFR can be understood as a variant of Follow-the-Regularized-Leader (FTRL)
applied independently at each information set. At each infoset, the acting
player minimizes cumulative regret using a linear loss with entropic
regularization (regret matching equals FTRL with the negative-entropy
regularizer).

This connection implies that CFR inherits FTRL's $O(1/\sqrt{T})$ regret
bound,
\[
  \frac{R^T}{T} \;\leq\; \frac{\Delta}{\sqrt{T}},
\]
where $\Delta$ is the range of utilities.

\subsection{CFR$^+$ as Optimistic FTRL}

The regret-flooring step in CFR$^+$ is equivalent to the ``optimistic''
update of Optimistic FTRL \citep{rakhlin2013optimistic}. Optimistic FTRL
achieves $O(1/T)$ regret when the loss sequence is predictable, which in
self-play corresponds to the strategy pair converging (each player's
strategy becomes more predictable as it approaches Nash). This provides a
theoretical mechanism for the faster empirical convergence of CFR$^+$.

\begin{proposition}
In self-play with alternating updates, CFR$^+$ with regret flooring is
equivalent to Optimistic FTRL with a specific predictive term. The $O(1/T)$
convergence observed empirically follows from the self-play structure
inducing low predictability error.
\end{proposition}

\subsection{Blackwell approachability}

The original CFR convergence proof \citep{zinkevich2007} is grounded in
Blackwell's approachability theorem \citep{blackwell1956analog}. The regret
minimization at each infoset corresponds to approaching the non-positive
orthant in regret space; CFR's convergence to Nash follows from the product
of per-infoset regret minimizers approaching a global regret minimum.
